\section*{Bab \ref{ch:b1:species-interactions} --- Interaksi Antarspesies}

\begin{solution}{exo:b1:species-interactions:1}
\omterm{def:b1:species-interactions:types}{Persaingan} $-/-$ (dua pohon ek memperebutkan cahaya); \omterm{def:b1:species-interactions:types}{pemangsaan} $+/-$
(lynx, terwelu); \omterm{def:b1:species-interactions:types}{parasitisme} $+/-$ (caplak, rusa); \omterm{def:b1:species-interactions:types}{mutualisme} $+/+$
(lebah, bunga); \omterm{def:b1:species-interactions:types}{komensalisme} $+/0$ (paku-pakuan pada kulit kayu);
amensalisme $-/0$ (rumput di bawah naungan sebatang pohon ek).
\end{solution}

\begin{solution}{exo:b1:species-interactions:2}
Dua spesies yang memakai sumber daya pembatas yang sama dengan cara
yang sama tak dapat berdampingan tanpa batas dalam sebuah \omterm{def:b1:organism-environment:organism}{lingkungan}
yang tetap. Kelima burung pengicaunya memakai makanan yang sama tetapi
di bagian pohon yang berbeda dan pada waktu yang berbeda: \omterm{def:b1:ecosystem-organization:niche}{relungnya}
berbeda, sehingga tak ada dua di antaranya yang ``dengan cara yang
sama''.
\end{solution}

\begin{solution}{exo:b1:species-interactions:3}
\omterm{def:b1:species-interactions:functional}{Waktu penanganan}: waktu untuk menangkap, memakan dan mencerna satu
mangsa, yang selama itu tak ada yang lain diambil. Laju serangan: luas
yang dicari tiap satuan waktu dikali peluang penangkapannya. Laju
maksimumnya $1/h = 2$ mangsa tiap jam.
\end{solution}

\begin{solution}{exo:b1:species-interactions:4}
Sebuah spesies yang penyingkirannya mengubah struktur \omterm{def:b1:ecosystem-organization:ecosystem}{komunitasnya}
tidak sebanding dengan kelimpahannya, biasanya dengan mengendalikan
sebuah pesaing yang mendominasi. Kelimpahan mengukur kehadiran, bukan
pengaruh; seekor pemangsa jarang yang memakan calon pemenangnya
memiliki sebuah pengaruh yang tak diperlihatkan jumlahnya.
\end{solution}

\begin{solution}{exo:b1:species-interactions:5}
\omterm{prop:b1:species-interactions:lotka}{Isoklin} 1: $N_1 + 0.5 N_2 = 500$ (potongannya 500 pada $N_1$, 1000 pada
$N_2$); \omterm{prop:b1:species-interactions:lotka}{isoklin} 2: $N_2 + 0.6 N_1 = 400$ (400 pada $N_2$, 667 pada $N_1$).
Keduanya berpotongan dengan tiap spesiesnya dibatasi lebih oleh dirinya ($\alpha = 0.5 <
K_1/K_2 = 1.25$; $\beta = 0.6 < K_2/K_1 = 0.8$): pendampingan mantap
pada $N_1 = 429$, $N_2 = 143$. Dengan $\alpha = 2$: \omterm{prop:b1:species-interactions:lotka}{isoklin} 1 memotong
pada 500 dan 250, keduanya di dalam potongan \omterm{prop:b1:species-interactions:lotka}{isoklin} 2 (667 dan 400):
spesies 2 selalu menang.
\end{solution}

\begin{solution}{exo:b1:species-interactions:6}
$f = aN/(1 + ahN)$: pada 10, $2/(1 + 0.5) = 1.33$ tiap jam; pada 100,
$20/(1 + 5) = 3.33$ tiap jam (maksimumnya $1/h = 4$). Setengah \omterm{def:b1:lipids:fattyacid}{jenuh}
pada $N = 1/(ah) = 20$ tiap meter persegi.
\end{solution}

\begin{solution}{exo:b1:species-interactions:7}
Dengan sebuah tanggapan jenis III fraksi mangsa yang dimakan naik
bersama kerapatannya pada kerapatan rendah: mangsa yang jarang diambil
secara sebanding lebih sedikit lalu dapat pulih. Dengan sebuah jenis
II fraksi yang dimakan paling tinggi pada kerapatan rendah
(pemangsanya tak pernah \omterm{def:b1:lipids:fattyacid}{jenuh} di situ), sehingga sebuah \omterm{def:b1:populations:population}{populasi}
mangsa yang kecil terdorong lebih jauh ke bawah.
\end{solution}

\begin{solution}{exo:b1:species-interactions:8}
\emph{P.~caudatum} memakai bakteri yang sama dalam kolom air yang sama
seperti \emph{P.~aurelia}, dan \emph{P.~aurelia} yang lebih cepat dan
lebih hemat menyeret makanannya di bawah apa yang diperlukannya.
\emph{P.~bursaria} makan di dasarnya pada sumber daya yang berbeda:
\omterm{def:b1:ecosystem-organization:niche}{relungnya}, bukan laju pertumbuhannya, yang memutuskan ---
\emph{P.~bursaria} tumbuh tidak lebih cepat daripada \emph{P.~caudatum}.
\end{solution}

\begin{solution}{exo:b1:species-interactions:9}
Pemangsanya berjumpa polanya lebih sering tanpa cedera, mempelajarinya
kurang baik, lalu menyerangnya lebih banyak: perlindungannya turun
bagi penirunya maupun modelnya. Karena itu kebugaran penirunya turun
seiring ia menjadi lazim, sebuah kebergantungan kekerapan yang negatif
yang menjaganya tetap jarang relatif terhadap modelnya.
\end{solution}

\begin{solution}{exo:b1:species-interactions:10}
Tanpa pemangsanya \omterm{prop:b1:species-interactions:lotka}{isoklin} kerang birunya terletak di luar \omterm{prop:b1:species-interactions:lotka}{isoklin} tiap
pesaingnya: ia memenangi batunya. \omterm{def:b1:species-interactions:types}{Pemangsaannya} menurunkan $K$ efektif
kerang birunya (\omterm{prop:b1:species-interactions:lotka}{isoklinnya} ditarik ke dalam) sampai ia memotong
\omterm{prop:b1:species-interactions:lotka}{isoklin} yang lain dari dalam: pesaingnya kini dibatasi lebih oleh
dirinya sendiri daripada oleh kerang birunya, lalu berdampingan
dengannya. Pemangsanya bekerja pada pesaing yang terkuat lalu membuat
ruang bagi sisanya.
\end{solution}

\begin{solution}{exo:b1:species-interactions:11}
Ketika fosfat membatasi pertumbuhannya, \qty{80}{\%} fosfatnya bernilai
jauh lebih daripada \qty{15}{\%} gulanya: \omterm{def:b1:species-interactions:types}{mutualisme}. Dalam sebuah
tanah yang kaya fosfat tumbuhannya akan menyerap fosfatnya sendiri dan
jamurnya memakan biaya \qty{15}{\%} tanpa guna: \omterm{def:b1:species-interactions:types}{parasitisme}.
Percobaannya: tumbuhkan tumbuhan dengan dan tanpa jamurnya pada
sebuah jangkau aras fosfat lalu ukurlah biomassanya; persekutuannya
bersifat mutualistik tempat tumbuhan bermikoriza tumbuh lebih besar,
parasitik tempat keduanya tumbuh lebih kecil.
\end{solution}

\begin{solution}{exo:b1:species-interactions:12}
Petak berpagar: petak berpagar di sepanjang sungainya yang tak dapat
dimasuki rusa elknya, dipasangkan dengan petak terbuka, diulang di
beberapa lembah dengan dan tanpa kelompok serigalanya, diukur
bertahun-tahun sebelum dan sesudah pemasukan kembalinya; catatlah
tinggi dan tutupan dedalunya, kerapatan rusa elknya dan waktu yang
dihabiskannya pada tiap petaknya (jejak, kamera), perburuan dan tebal
saljunya. Riamnya meramalkan dedalunya pulih pada petak terbuka hanya
tempat serigalanya ada, pada laju yang sama seperti di dalam petak
berpagarnya; bila dedalunya pulih sama saja dengan dan tanpa
serigalanya, atau mengikuti salju dan perburuannya alih-alih itu,
riamnya tersangkal atau terbagi dengan faktor itu.
\end{solution}

\begin{solution}{pb:b1:species-interactions:1}
\textbf{1.} A: $N_A + 1.4 N_B = 200$, potongannya 200 ($N_A$) dan 143
($N_B$). B: $N_B + 0.9 N_A = 130$, potongannya 130 ($N_B$) dan 144
($N_A$).
\textbf{2.} Potongannya (200 lawan 144 pada sumbu $N_A$, 143 lawan
130 pada sumbu $N_B$) menaruh \omterm{prop:b1:species-interactions:lotka}{isoklin} A seluruhnya di luar \omterm{prop:b1:species-interactions:lotka}{isoklin} B:
menyelesaikan $N_A = 200 - 1.4 N_B$ dengan $N_B = 130 - 0.9 N_A$ memberi $N_A = -69$,
sehingga garisnya bertemu di luar kuadran positifnya. A menang dari
titik awal mana pun: di mana pun B masih dapat tumbuh, A tumbuh juga
lalu bertahan lebih lama.
\textbf{3.} $T_A = 0.87$ hari, $T_B = 1.26$ hari: A, yang tumbuh lebih
cepat.
\textbf{4.} A: $N_A + 0.3 N_C = 200$ (200, 667); C: $N_C + 0.3 N_A =
100$ (100, 333). Tiap spesiesnya membatasi dirinya lebih daripada yang
lain ($0.3 < 200/100$ dan $0.3 < 100/200 = 0.5$): pendampingan mantap.
\textbf{5.} Sebuah hutan tidak tetap (musim dan gangguannya menyetel
ulang perlombaannya sebelum ia diputuskan) dan memiliki banyak sumber
daya serta tempat, sehingga spesiesnya membagi \omterm{def:b1:ecosystem-organization:niche}{relung} alih-alih
berbagi satu.
\textbf{6.} $N_A = 200 - 0.3 N_C$, $N_C = 100 - 0.3 N_A = 100 - 60 +
0.09 N_C$, sehingga $N_C = 44$ dan $N_A = 187$.
\textbf{7.} Naik lalu mendatar menuju sekitar 10: jenis II.
\textbf{8.} $(1/N, 1/f)$: $(0.2, 0.5)$, $(0.1, 0.303)$, $(0.05, 0.2)$,
$(0.025, 0.149)$, $(0.0125, 0.125)$: sebuah garis lurus.
\textbf{9.} Kemiringannya $(0.5 - 0.125)/(0.2 - 0.0125) = 2.0 = 1/a$, sehingga $a =
\qty{0.5}{m^2/d}$; potongannya $0.5 - 2\times 0.2 = 0.1$ hari $= h$.
\textbf{10.} Maksimumnya $1/h = 10$ mangsa tiap hari; setengah \omterm{def:b1:lipids:fattyacid}{jenuhnya} $N =
1/(ah) = 20$ tiap meter persegi.
\textbf{11.} Fraksi penanganannya $= f h = f/10$: lebih dari separuh ketika $f
> 5$, yakni pada 40 dan 80 tiap meter persegi (5,0 pada 20 tepat separuh).
\textbf{12.} Pertumbuhan mangsanya $0.1 N (1 - N/200)$ setara dengan
pemakaiannya $0.5N/(1 + 0.05N)$: $0.1(1 - N/200) = 0.5/(1 + 0.05 N)$, yakni
$(1 - N/200)(1 + 0.05N) = 5$: $1 + 0.05N - N/200 - N^2/4000 = 5$, sehingga
$N^2 - 180 N + 16\,000 = 0$, $N = 90 \pm \sqrt{8100 - 16\,000}$: tak ada
akar nyata --- satu kumbang tiap meter persegi memakan lebih daripada
yang pernah sanggup dihasilkan mangsanya (pertumbuhan maksimumnya $rK/4 = 5$
tiap hari, pada $N = 100$, ketika kumbangnya memakan 8,3) lalu
menyeretnya sampai punah. Tak ada kesetimbangan; mangsanya bertahan
hanya bila kumbangnya lebih sedikit atau mangsanya memiliki
persembunyian.
\textbf{13.} \omterm{def:b1:species-interactions:functional}{Waktu penanganannya} terparuh bagi mangsa keduanya ($h' = 0.05$),
sehingga pada kejenuhannya kumbangnya memakan sampai 20 di antaranya
sehari: asupan totalnya pada kerapatan tinggi kira-kira berlipat dua
bila ia beralih ke mangsa yang lebih cepat.
\textbf{14.} $H_{15} = \ln 15 = 2.71$. Delapan spesies: $-(0.85\ln 0.85
+ 7\times 0.0214\ln 0.0214) = 0.138 + 0.576 = 0.71$.
\textbf{15.} Sebesar faktor $2.71/0.71 = 3.8$.
\textbf{16.} Tanda $-$ bintang lautnya pada kerang birunya menjaga
tanda $-$ kerang birunya pada pesaingnya tetap lemah; singkirkanlah
yang pertama dan kerang birunya mencapai kerapatan ketika
\omterm{def:b1:species-interactions:types}{persaingannya} menyisihkan sisanya.
\textbf{17.} Biomassa mengukur berapa banyak alir energi yang dipegang
sebuah spesies, bukan apa yang diperbuat interaksinya; pengaruh sebuah
\omterm{def:b1:species-interactions:keystone}{spesies kunci} berjalan lewat spesies yang dikendalikannya, yang
biomassanya besar.
\textbf{18.} Menyingkirkan bintang lautnya akan melepaskan sebuah
spesies jarang yang tak sanggup mengambil alih batunya: \omterm{def:b1:ecosystem-organization:ecosystem}{komunitas} yang
didominasi kerang birunya akan sedikit berubah, dan bintang lautnya
tidak akan menjadi sebuah \omterm{def:b1:species-interactions:keystone}{spesies kunci}.
\textbf{19.} Sangkar yang menghalangi bintang lautnya dari petak
kecil, dengan sangkar terbuka sebagai kendali bagi pengaruh sangkarnya
sendiri: bila kerang birunya mengambil alih hanya dalam sangkar yang
tertutup, \omterm{def:b1:species-interactions:types}{pemangsaan} atas kerang birunya adalah mekanismenya; sangkar
dengan bintang lautnya hadir tetapi kerang birunya disingkirkan dengan
tangan akan menguji apakah mangsa bintang lautnya yang lain berperan.
\textbf{20.} Perolehannya $\qty{40}{mg}$, biayanya $8\times 0.5 = \qty{4}{mg}$: bersihnya
\qty{36}{mg} $=$ \qty{0.58}{kJ}.
\textbf{21.} \qty{40}{\micro g} dari \qty{4}{mg}: \qty{1}{\%}
fotosintat harian bunganya membeli sekitar tiga pemindahan serbuk
sari, yaitu pembiakan tumbuhannya.
\textbf{22.} Masing-masing mengambil sesuatu dari yang lain, tetapi
masing-masing memperoleh lebih daripada yang hilang darinya: tandanya
adalah tanda pengaruh bersihnya, bukan tanda eksploitasinya.
\textbf{23.} Perampoknya $+$, tumbuhannya $-$ (nektarnya diambil, tanpa
penyerbukan): sebuah \omterm{def:b1:species-interactions:types}{parasitisme} atas \omterm{def:b1:species-interactions:types}{mutualismenya}; lebahnya menemukan
bunga yang terkosongkan lalu memperoleh lebih sedikit, sebuah $-$ tak
langsung.
\textbf{24.} Tumbuhan $\leftrightarrow$ lebah $+/+$; perampok $\to$ tumbuhan $-$,
perampok $\to$ lebah $-$; burung $\to$ perampok $-$. Dua tanda kurang:
burungnya memiliki sebuah $+$ tak langsung pada tumbuhannya dan pada
lebahnya.
\textbf{25.} $a = \qty{0.5}{m^2/d}$, $h = \qty{0.1}{d}$ (\qty{2.4}{h}),
maksimumnya \num{10} mangsa tiap hari.
\end{solution}
